% Person B handoff. Include in the private eandd.tex or paste each section in its place.
% Requires the booktabs package (\toprule, \midrule, \bottomrule).
\section{Data and evaluation protocol}
\label{sec:setting-online}
\paragraph{Data.} The data are a binary response matrix $Y\in\{0,1\}^{N\times M}$ with $N$ subjects (rows, index $i$) and $M$ items (columns, index $j$): $Y_{ij}=1$ if subject $i$ answers item $j$ correctly. Most entries are missing, since a subject answers only the items of the benchmarks it was evaluated on. Each item $j$ belongs to the dataset $d(j)$ of one of $D=17$ benchmarks, so a benchmark is a block of columns, and a \emph{task} is one row restricted to one block: subject $i$ on benchmark $d$.

Each item $j$ has an embedding $x_j\in\mathbb R^{1024}$ computed once from its text by the frozen BGE-M3 encoder~\citep{chen2024bge} with a context of 8,192 tokens, and normalised to unit length. We use BGE-M3 because the corpus is multilingual (it includes Chinese, Russian and several African languages) and some items are long, which English-only and 512-token encoders handle poorly.

\paragraph{Online simulation and splits.} The competition protocol is \emph{online}: for each task, labels are revealed one at a time and the predictor updates from them. Whatever the predictor knows before the first label is \emph{offline}, that is, fitted on the training benchmarks beforehand. Our \emph{online simulation} runs this protocol on recorded responses, with disjoint column blocks for training, validation, and test (11, 3, and 3 benchmarks): priors are fitted offline on training, and validation selects settings. The validation and test benchmarks are \emph{out of distribution} (OOD): their responses are excluded from offline fitting.

A \emph{competition-shaped draw} samples one subject per benchmark, so a predictor cannot borrow labels from other subjects on the same benchmark; this simulates the competition protocol on public data. As a separate \emph{in-distribution} (ID) evaluation, we hold out rows (subjects) within the training blocks, so its benchmarks and items have been seen during fitting, and retain the questions covered by the fit subjects. The offline comparator reuses its fitted benchmark rates and item difficulties where available, and its online update changes only the subject ability.

\paragraph{Budgets and scoring.} For each subject--benchmark pair we freeze a split of distinct items into acquisition and evaluation pools, assigning the extra item to acquisition when the count is odd. A frozen random acquisition order supplies nested prefixes of 0, 1, 3, 7, 15, and 31 labels. At each budget the predictor conditions on the allowed prefix and predicts the same evaluation items.

Evaluation outcomes are used only to score those predictions. We average Brier scores over query items, then average repeated appearances within each task and give tasks equal weight. Brier ALC combines the six budgets with weights $(0.1,0.2,0.2,0.2,0.2,0.1)$. Paired intervals resample whole tasks with their matched predictions.

\section{Probes}
\label{sec:probes-text}
The predictors above ignore what an item says. A probe predicts an item's relative difficulty from its embedding before any response from that benchmark is available. We compare ridge regression on the embeddings, a probe using log character count, and a random score. We add ridge regression and a small neural network (MLP) that fit either the item pass rate or the negative fitted Rasch difficulty (easiness), each centred within its benchmark. The fitted probes use a fixed sample of 8,462 items from the 11 training benchmarks, at most 800 per benchmark, and choose their regularisation by five folds grouped by training benchmark. Appendix~\ref{app:probes} gives the training details.

We score a probe on each benchmark $d$ by two numbers and average each over benchmarks. Let $y_j$ be the observed pass rate of item $j$ across recorded subjects and $\hat y_j$ its predicted easiness, so higher values mean easier items. The \emph{Spearman correlation} $\rho_d=\operatorname{corr}\big(\operatorname{rank}(\hat y_j),\operatorname{rank}(y_j)\big)$ compares the two orderings of the items: it is $1$ if the probe orders them perfectly, $0$ if its order is unrelated to the truth, and $-1$ if it is reversed. The \emph{pairwise accuracy} $A_d=\frac{1}{|P_d|}\sum_{(j,k)\in P_d}s_{jk}$ is the share of item pairs the probe orders correctly, where $P_d$ is the set of pairs with $y_j\neq y_k$ (observed ties are excluded), $s_{jk}=1$ if $\hat y_j-\hat y_k$ has the same sign as $y_j-y_k$, $s_{jk}=\tfrac12$ if the probe predicts a tie, and $s_{jk}=0$ otherwise. A random score has $A_d=0.5$ on average.

To test whether a probe helps prediction, we give the item deviations of Rasch a prior mean taken from it,
\begin{equation}\label{eq:probe-prior}
\delta_j\sim\mathcal N\big(\gamma\,z_j,\ \sigma_\delta^2\big),
\end{equation}
where $z_j$ is the probe's predicted difficulty of item $j$, centred within its benchmark (the benchmark rate $\mu_d$ absorbs the average), and $\gamma$ is a scale. Setting $\gamma=0$ gives the fitted Rasch. At prediction time the probes use only item features. Training outcomes supply supervision, and validation and test outcomes are used only for evaluation and for choosing $\gamma$ on validation.

\section{Item features}
\label{sec:itemfeat-text}
\begin{table}[H]
\centering
\caption{Item-difficulty transfer: within-benchmark Spearman correlation and pairwise accuracy between predicted easiness and observed pass rate, averaged across benchmarks. ID: items held out within the 11 training benchmarks; OOD validation and OOD test: three unseen benchmarks each. The additional pass-rate and Rasch-target probes use benchmark-grouped selection and were not evaluated on ID (--).}
\label{tab:item-difficulty}
\small\setlength{\tabcolsep}{4pt}
\input{eandd_probe_table}
\end{table}

Embeddings show a clear signal on known benchmarks. On the six unseen benchmarks the two MLP probes pass the criterion we fixed in advance, with benchmark-macro Spearman $0.090$ and $0.103$ against $0.026$ for a random baseline and a positive correlation on 5 of 6 benchmarks; the ridge probes do not ($0.065$ and $0.047$). On the three test benchmarks alone, the pass-rate MLP reaches $0.109$, while the Rasch-target MLP reaches $0.059$, below the sampled random control's $0.076$. The aggregate rank criterion therefore does not establish a consistent test-only advantage.

The rank signal yields only a small gain at the validation-selected scale. In Table~\ref{tab:probe-prior}, every $\gamma\ge0.25$ is worse than $\gamma=0$ on both splits. The scale selected on validation ($\gamma=0.05$, Rasch-target MLP) improves Brier ALC by about $0.0004$ on each split in its matched replay; this comparison alone does not establish a reliable improvement. This tests these two probes only; a stronger difficulty predictor, for instance one trained through an item response likelihood on more items, could do better.

\begin{table}[H]
\centering
\caption{Predicted item difficulty as the prior mean of the Rasch item deviation, Equation~\eqref{eq:probe-prior}. Cells are $\Delta_{m-\mathrm{J}}$ of the fitted Rasch with scale $\gamma$, competition-shaped draws; $\gamma=0$ is the fitted Rasch without a prior mean. Bold: lowest in each column.}
\label{tab:probe-prior}
\input{eandd_r5}
\end{table}

\section{Limitations}
\label{sec:limitations}
The filtered public data covers only three test benchmarks in the online simulation. Multiple draws
and acquisition orders measure variation within this population; they do not
supply independent tests of transfer to new benchmarks. Task bootstrap
intervals likewise do not capture uncertainty from sampling new benchmarks.
Results apply to the eligible response data and random label acquisition,
and need not carry over to other benchmark mixes or selection policies.
The 17-benchmark sensitivity analysis reuses held-out-subject tasks for
the 11 original training benchmarks. The test triples overlap, so their
distribution describes sensitivity within this dataset rather than
independent estimates of transfer.
The ID comparison has item response history that is unavailable under
benchmark cold start. Observed item difficulty also depends on which subjects
answered the item. The leave-one-benchmark-out association uses realized
outcomes to explain when the fitted prior helps; it is not a way to select a
prior before observing a new benchmark.

\section{Disclosure}
\label{sec:disclosure-rerun}
\textbf{Training data.} All fitted quantities (the offline Rasch fit, the Beta-Binomial
empirical-Bayes prior, the item probes) use only the responses of the 11 training benchmarks in
the competition's measurement database after the filtering described in Section~\ref{sec:setting-online};
no other response data is used. \textbf{Pretrained models.} Item text is embedded with the frozen
BGE-M3 model; nothing else is pretrained or fine-tuned. \textbf{LLM coding assistants.} Claude Code
(Anthropic) assisted with writing code, analysis scripts and drafting text; every reported number was
produced by running the code in the repository, and the authors checked the numbers.
OpenAI Codex assisted with the reviewer analyses, figures, and text revisions.
